\abstract{РЕФЕРАТ}

Объем работы равен \formbytotal{lastpage}{страниц}{е}{ам}{ам}. Работа содержит \formbytotal{figurecnt}{иллюстраци}{ю}{и}{й}, \formbytotal{tablecnt}{таблиц}{у}{ы}{}, \arabic{bibcount} библиографических источников. Количество приложений -- 2. Графический материал представлен в приложении А. Фрагменты исходного кода представлены в приложении Б.

Перечень ключевых слов: регулярные выражения, метасимволы, шаблон, библиотека, компиляция шаблона, конечный автомат, рекурсивный спуск, поиск с возвратами, группы захвата, обратные ссылки, символьные классы, квантификатор, абстрактное синтаксическое дерево, C API, Qt, текстовый редактор, графический интерфейс, поиск по директории, UTF-8, язык Си, C++, кроссплатформенность.

Объектом разработки являются библиотека регулярных выражений на языке Си с переносимым C API и десктопное поисковое приложение с графическим интерфейсом на основе фреймворка Qt, интегрирующее разработанную библиотеку.

Целью выпускной квалификационной работы является проектирование и реализация программно-информационной поисковой системы с графическим интерфейсом на основе библиотеки регулярных выражений собственной разработки.

Библиотека регулярных выражений реализована на языке Си методом рекурсивного нисходящего парсера: шаблон компилируется в абстрактное синтаксическое дерево непосредственно в процессе однопроходного разбора без построения промежуточных представлений, а сопоставление с текстом выполняется алгоритмом поиска с возвратами над полученным деревом.

Разработанная библиотека регулярных выражений может быть использована как самостоятельный компонент в любом проекте на языке Си или C++, требующем поиска, валидации или обработки текстовых данных. Поисковое приложение может применяться разработчиками и системными администраторами для интерактивного поиска и редактирования текстовых файлов и журналов с использованием регулярных выражений.



\selectlanguage{english}
\abstract{ABSTRACT}

The volume of work is \formbytotal{lastpage}{page}{}{s}{s}. The work contains \formbytotal{figurecnt}{illustration}{}{s}{s}, \formbytotal{tablecnt}{table}{}{s}{s}, \arabic{bibcount} bibliographic sources. The number of appendices is 2. Graphic material is presented in appendix A. Source code fragments are presented in appendix B.

Keywords: regular expressions, wildcards, pattern, library, pattern compilation, finite automaton, recursive descent, backtracking, capture groups, backreferences, character classes, quantifier, abstract syntax tree, C API, Qt, text editor, graphical user interface, directory search, UTF-8, C language, C++, cross-platform.

The object of development is a regular expression library written in C with a portable C API, and a desktop search application with a graphical user interface based on the Qt framework, integrating the developed library.

The purpose of the graduation thesis is to design and implement a software information retrieval system with a graphical user interface based on a custom-developed regular expression library.

The regular expression library is implemented in C using recursive descent parsing: the pattern is compiled into an abstract syntax tree directly during a single-pass parse without constructing intermediate representations, and text matching is performed by a backtracking search algorithm over the resulting tree.

The developed regular expression library can be used as a standalone component in any C or C++ project requiring text search, validation, or processing. The search application can be used by developers and system administrators for interactive search and editing of text files and logs using regular expressions.

\selectlanguage{russian}